% ID: SEL-D0048-OA-E01
% ESTADO: cotejada. Segunda pasada visual y borrador compilado revisados.
% Fuente: Selectividad/sel_2013_matematicas_aplicadas/junio/reserva a/EXAMEN 5 -ANDALUCÍA.pdf
% SHA-256: 76b6f6a138ec0917a2e3b3a2925bccca739af3363e419603b6814d835860441e
% Página: 1; opción A; ejercicio original 1.

\begin{enumerate}[label=\alph*)]
\item \textbf{(1 punto)} Plantee, sin resolver, el siguiente problema:

``Un barco puede transportar vehículos de dos tipos: coches y motos. Las condiciones de la nave obligan a que el número de motos no pueda ser inferior a la cuarta parte del de coches ni superior a su doble; además, la suma del número de motos más el doble del número de coches no puede ser mayor que 100. \textquestiondown{}Cuántos vehículos, como máximo, puede transportar este barco?''
\item \textbf{(1.5 puntos)} Dado el recinto limitado por las inecuaciones
\[
y\geq30,\qquad3x-y\geq150,\qquad6x+7y\leq840,
\]
halle en qué puntos de ese recinto la función $F(x,y)=6x-2y$ alcanza su valor mínimo.
\end{enumerate}
